\subsection*{Exterior-tail certification of the A.11 axial cone component}
\label{subsec:A11-actual-Ns-tail}
On the two corrected supports of (A.11), the source has $U=M=J=0$.
After the genuine amplitude equation imposes $S(\infty,\eta)=0$, the source
moment equation (4.15) gives
\begin{equation}\label{eq:A11-tail-SP}
 S(X,\eta)=\frac12\int_X^\infty E(x,\eta)^2\,dx,\qquad
 \Pi(X,\eta)=-\frac12\int_X^\infty\frac{E(x,\eta)^2}{x}\,dx.
\end{equation}
The pressure formula uses the actual condition $\Pi(\infty,\eta)=0$
and the A.11 pressure-neutral moment correction. In particular
the $S$ identity is not valid before the $S(\infty)=0$ condition.

\begin{lemma}[Exact tail form of the axial stress source]
Define $D_s=S/(XE^2)$, $P_s=-\Pi/E^2$, $g=\partial_\eta\log E$,
$d=1-\eta^2$, $A=1/2+h$. Formula (4.16) becomes
\begin{equation}\label{eq:A11-Ns-tail-exact}
 \frac{N_s}{E}=E\left\{4\eta(hD_s-AP_s)
 -d[2g(D_s-P_s)+\partial_\eta D_s-\partial_\eta P_s]\right\}.
\end{equation}
If $\ell=\partial_{\log X}\log(\sqrt X E)\le0$ at all later radii
and $|\partial_\eta\log E|\le\varepsilon$ there, then
\begin{equation}\label{eq:A11-Ns-tail-inequality}
0\le P_s\le\frac12,\qquad
\frac{|N_s|}{E}\le E[4hD_s+2A+2\varepsilon(D_s+1/2)].
\end{equation}
\end{lemma}
\begin{proof}
For $t=\log(x/X)$ write $R(t)=\exp(2\int_0^t\ell(Xe^v)\,dv)$.
Then $D_s=\frac12\int_0^\infty R\,dt$ and
$P_s=\frac12\int_0^\infty e^{-t}R\,dt$. Since $R\le1$, $P_s\le1/2$.
The original tail integrals give
$|S_\eta|/X\le2\varepsilon E^2D_s$ and
$|\Pi_\eta|\le2\varepsilon E^2P_s$. Apply (4.16).
The exact displayed expression follows by differentiation.
\end{proof}

\begin{theorem}[Source-stage sufficient bound for both A.11 supports]
\label{thm:A11-Ns-realized}
Retain the corrected two-bump bounds, $T_f=512$, $T_d>10$,
$P_*>e^{T_d}$ and the post-pulse nonpositive slope and terminal
release bounds used above. Choose
\[
0<\lambda\le\min\left\{
\frac{10^{-6}}{T_d+513},\frac1{10(1+\log P_*)},
\frac1{20}\right\},\qquad0<h\le\lambda/100.
\]
Once the actual source moment matching has imposed $S(\infty)=0$,
the resulting corrected supports satisfy
\begin{equation}\label{eq:A11-Ns-final}
\sup_{\mathrm{A.11\ supports},\,|\eta|\le1}|N_s|/E<\lambda/2.
\end{equation}
Together with $Q_s\ge\lambda/2$, the corrected slope range
$-3\lambda/2\le\ell\le-\lambda/2$, and a choice
$X\ge9/\lambda$, the full strict cone (4.22) follows on these
two supports, with the previous normalized margin $7/20$.
\end{theorem}
\begin{proof}
Past the angular interpolation the unedited source $E$ is
$\eta$-independent. Only the two corrected bumps retain
$\eta$ dependence. The bounds
$\|c_\eta\|_\infty<12\lambda^{28}$,
$\|\beta\|_\infty\le640/3$ and
$|\sum c_i\beta_i|<1/2$ give
$|\partial_\eta\log E|<5120\lambda^{28}=:\varepsilon$
at every radius beyond the first support.
The prescribed later nonpositive slope stages imply
\[
D_s\le B_D:=T_f/2+15\log(1/\lambda)+3+h^7/3
\]
on either support: allow half the remaining interpolation/plateau
length, the two release transitions, the $-1$ hold and its
exponentially suppressed terminal tail.

The first unit transition enlarges $E$ by at most $e^{1/10}$.
The angular flattening has integrated positive contribution at
most $\log2<1$, and the relative correction is less than $3/2$.
The corrected supports lie after a length-$13/\lambda$ pulse
and within a three-unit endpoint reserve. Dropping all other
decay gives
\[
 E\le P_*\exp(4-13/(2\lambda)).
\]
Put $t=1/\lambda\ge20$; since
$t\ge10(1+\log P_*)$ and $\log(4t)\le t$,
\[
\log(4tP_*e^{4-13t/2})\le3-27t/5<0,
\]
hence $E<\lambda/4$.

Also $\lambda T_f\le10^{-6}$, $h\le\lambda/100$,
$\lambda\log(1/\lambda)\le1$ and $\lambda\le1/20$ give
$4hB_D<61/100$, $2A\le1001/1000$, and
$2\varepsilon(B_D+1/2)<1/100$ (use
$\lambda^{27}\le20^{-27}$). Thus the square bracket in
\eqref{eq:A11-Ns-tail-inequality} is strictly less than $2$,
so $|N_s|/E<\lambda/2$. The previously proved
$Q_s$ and slope bounds plus $X\ge9/\lambda$ imply the
source cone inequality at these supports.
\end{proof}

This is an all-parameter \emph{conditional source-stage theorem},
not an enclosure of one independently frozen complete exterior
schedule. Compatibility with every other source-stage stress-cone
margin, pulse $o(1)$ constants, terminal matching, Appendix B
and the unforced coefficient hierarchy still requires separate
verification. The proposed unforced Theorem~1.1 remains open.
